% Section content template for individual Educative lessons
% This file contains the actual content for a specific section
% Generated from Educative API section content

% Section: Key Insights: The Hotel Management System
% Section ID: 5901566010589184
% Section Slug: key-insights-the-hotel-management-system
% Generated: 2025-12-12T16:23:52.456625

Well done on completing the hotel management system case study! This lesson takes you further by highlighting frequent design mistakes, sharing actionable ways to refine your design thinking, and equipping you with insights to handle tricky interview questions. You 'll also discover additional case studies to deepen your grasp of practical system design concepts.

\subsection{Common mistakes}\label{e2-gHdeBFlSE6eHtXbHkc}
\\
Designing the Hotel Management System can be tricky, and candidates often fall into similar traps that weaken their solutions.

\begin{itemize}
\item
\protect\phantomsection\label{qS63wCpzUFhvGF4MPUCle}
\textbf{Overlooking concurrency in bookings:} Many candidates forget to handle double-booking scenarios. For example, two guests could book the same room if overlapping bookings aren't properly synchronized. Design the \texttt{RoomBooking} logic so that checking a room's availability and marking it as booked happen together as one single step. This ensures that no other booking can claim the same room simultaneously.
\item
\protect\phantomsection\label{uII4V45ph5YJL87NMDj8a}
\textbf{Unclear separation of responsibilities:} It's common to mix roles, such as letting the \texttt{Guest} class handle room status updates instead of the \texttt{Receptionist} or \texttt{System} classes. Keep each class focused on its responsibility. For example, the \texttt{Guest} only requests actions; \texttt{receptionists} and \texttt{Systems} update room status and issue keys.
\item
\protect\phantomsection\label{FzDTQAjG0MAf5Qw8q8zb7}
\textbf{Ignoring edge cases for cancellations:} Candidates often neglect refund rules, such as the requirement that full refunds be applied only if a booking is canceled at least 24 hours before check-in. Clearly define cancellation rules in the \texttt{RoomBooking} logic and update \texttt{Invoice} and \texttt{PaymentStatus} accordingly to handle partial or no refunds.
\item
\protect\phantomsection\label{J-stZlqY2kZXEa-OmxGp5}
\textbf{Missing proper pricing strategy handling:} Some candidates hard-code room prices and forget that different rooms may require flexible pricing models (e.g., seasonal rates or special discounts). Use the \texttt{Strategy} pattern to separate and switch pricing algorithms easily, so each room type can apply its pricing logic when needed.
\item
\protect\phantomsection\label{FZrcVb7vateD9HRnR9XUm}
\textbf{Overlooking hotel hierarchy relationships:} Many designs treat the hotel as a single unit and don't model how multiple branches, each with its rooms and bookings, fit into the bigger structure. Use the \texttt{Hotel} and \texttt{HotelBranch} classes to represent a clear hierarchy supporting multiple branches, each managing its inventory independently.
\end{itemize}

\subsection{Key interview tips}\label{UCUsXYsxYIqUCYUSynwVo}
\\
Small changes in how you present your design can make a big difference. Here are tips to communicate your Hotel Management System design effectively:

\begin{table}[h!]
\centering
\begin{tabular}{|p{6cm}|p{9cm}|}
\hline
\textbf{Tip} & \textbf{Explanation} \\
\hline
Prioritize user flows. & 
Before creating class diagrams, outline the most critical user flows (e.g., booking a room, checking in). This demonstrates a strong understanding of requirements and sets the context for the design, making it easier for interviewers to follow the thought process. \\
\hline
Explain account types. & 
Clarify how Guest, Receptionist, and Housekeeper extend the base Person class to keep responsibilities clear and modular. \\
\hline
Discuss notifications. & 
Explain how notifications (SMS and Email) are handled through the Notification abstract class to keep guests informed automatically. \\
\hline
Highlight system extensibility. & 
Describe how adding new room styles, services, or payment types does not break existing classes due to the use of abstract classes and clear interfaces. \\
\hline
Mention fail-safe flows. & 
Discuss how the design handles failures, such as payment failure during booking or abandoned booking scenarios. \\
\hline
\end{tabular}
\caption{Design and Explanation Tips for a Hotel Management System}
\label{tab:hotel_design_tips}
\end{table}


\subsection{Test your understanding}\label{7_dnqls0DZO2dN-d8JZ5n}
\\
Test your knowledge with these questions and confirm you understand how to design and explain this system well.

\subsection*{The Hotel Management System}
\vspace{0.3cm}
\noindent
\textbf{Question 1:} Which relationship best describes how \texttt{RoomBooking} relates to \texttt{Invoice}?
\\[0.2cm]
\begin{itemize}
 \item Association
 \item Aggregation
 \item \textbf{[CORRECT]} Composition
\\[0.1cm]
\textit{Explanation:} 
The \texttt{RoomBooking} class is composed of the \texttt{Invoice} class, i.e., it cannot exist independently
 \item Generalization
\end{itemize}
\vspace{0.3cm}
\noindent
\textbf{Question 2:} Which principle does extending the \texttt{Person} class into subclasses \texttt{Guest}, \texttt{Receptionist}, and \texttt{Housekeeper} best represent?
\\[0.2cm]
\begin{itemize}
 \item Open/Closed principle
 \item \textbf{[CORRECT]} Single Responsibility principle
\\[0.1cm]
\textit{Explanation:} 
Each subclass has clear, separate duties
 \item Liskov Substitution principle
 \item Dependency Inversion principle
\end{itemize}
\vspace{0.3cm}
\noindent
\textbf{Question 3:} Which design element helps manage different room services like amenities and kitchen orders?
\\[0.2cm]
\begin{itemize}
 \item The \texttt{Room} class directly handles them
 \item \textbf{[CORRECT]} A \texttt{Service} abstract class with child classes
\\[0.1cm]
\textit{Explanation:} 
A \texttt{Service} abstract class with subclasses \texttt{Amenity}, \texttt{RoomService}, and \texttt{KitchenService} manages extra services.
 \item The \texttt{Catalog} class
 \item The \texttt{Account} class
\end{itemize}
\vspace{0.3cm}
\noindent
\textbf{Question 4:} Why does the design use a Strategy pattern for room pricing instead of putting price logic inside the \texttt{Room} class?
\\[0.2cm]
\begin{itemize}
 \item To increase inheritance hierarchy depth.
 \item To keep the \texttt{Room} class open for modification.
 \item \textbf{[CORRECT]} To follow the Open/Closed principle and allow new pricing algorithms without changing \texttt{Room}.
\\[0.1cm]
\textit{Explanation:} 
Using Strategy keeps the \texttt{Room} class closed for modification but open for extension with new pricing logic.
 \item To reduce the number of classes.
\end{itemize}
\vspace{0.3cm}
\noindent
\textbf{Question 5:} If the \texttt{Receptionist} checks in a guest but forgets to issue a \texttt{RoomKey}, which part of the design could detect or prevent this inconsistency?
\\[0.2cm]
\begin{itemize}
 \item \textbf{[CORRECT]} The \texttt{Housekeeper} class should check the key status.
 \item The \texttt{Hotel} Singleton should auto-issue keys.
 \item The \texttt{Search} interface should reject the booking.
 \item The \texttt{RoomBooking} process should include an explicit check-key step.
\\[0.1cm]
\textit{Explanation:} 
The \texttt{RoomBooking} flow should enforce a step to check that a key is issued when check-in is done. Otherwise , the room can't be securely accessed.
\end{itemize}

\subsection{Related case studies}\label{K3mySFDAVTtf0SgiSlcQb}
\\
Explore these related case studies to deepen your understanding of similar design challenges and reusable design ideas.

\begin{itemize}
\item
\textbf{Designing a car rental system:} Similar to hotel room booking, this case study tests your approach to handling availability, overlapping bookings, dynamic pricing, and branch-specific inventory.
\item
\textbf{Designing an airline management system:} Reinforces concepts like seat reservation conflicts, cancellation policies, and notification flows, just like hotel room booking.
\item
\textbf{Designing a restaurant management system:} This will help you practice handling table reservations, service requests, staff roles, and payments, closely mirroring hotel dining or room service add-ons.
\item
\textbf{Designing a movie ticket booking system:} Strengthens your thinking about time-slot-based booking, seat selection, preventing double bookings, and handling online payments.
\item
\protect\phantomsection\label{Ud9k7PcEfhGZCIp5SzLqx}
\textbf{Designing the Amazon online shopping system:} Builds on the idea of handling multiple user accounts, dynamic pricing strategies (discounts, special offers), and payment workflows similar to hotel service add-ons and billing.
\end{itemize}

% End of section content